\textbf{Background:} Connecting serum inflammation with clinical severity requires a practical calculation that remains interpretable across treatment and follow-up. We evaluated a common Eczema Area and Severity Index (EASI) equation within a Chinese treat-to-target cohort.

\textbf{Methods:} We examined clinical outcomes and all 23 serum markers under dupilumab and two initial abrocitinib doses, adjusting treatment comparisons for measured baseline differences. A model-development procedure using nominal-visit-matched affected body surface area (BSA), CCL17, CCL26 and IL-5 was developed with baseline/week-12 observations using nested patient-level validation and applied at week 24. Individual estimates, within-person changes and EASI75 identification were evaluated separately.

\textbf{Results:} Clinical comparisons included 163 patients; EASI estimation included 316 visits from 160 patients. Clinical control improved alongside treatment-dependent molecular changes. The selected four-input equation required no treatment label or previous score; performance estimates describe the nested development procedure. Mean absolute error (MAE) at baseline/weeks 12/24 was 3.19/1.30/1.19, with Spearman correlations of 0.918/0.927/0.904. Follow-up errors were within two points in 84.3\% and 82.0\% of predictions. MAEs for within-person changes over baseline--week 12, baseline--week 24 and weeks 12--24 were 3.87, 3.47 and 1.49. EASI75 sensitivity at weeks 12/24 was 76.5\%/86.1\%, with specificity of 98.6\%/96.2\%.

\textbf{Conclusions:} The BSA--serum development procedure gave close follow-up EASI estimates and tracked changes within this cohort. Molecular added accuracy was not established, and baseline severe disease remained the main accuracy boundary. These results support independent evaluation of this restricted-input research measure across the severity range.
